\scanpage{96}
\mspage{11''}
% archive sheet number: 95

\uline{Remarque} \unsure{2.7 bis|2.7}. Il semble plausible qu'on doive pouvoir montrer, \deleted{\unsure{à}} \inserted{lorsqu'on dispose de} la résolution des singularités, que les faisceaux $\underline{\unsure{\Phi|\varphi}}_\nu^{(q)}$ \inserted{(ou les $\underline\Psi_\nu^{(q)}$, c'est pareil)}, \deleted{\unsure{sera}} sont constructibles si $\ell\ne\deleted{\unsure{\mathrm{car}}}p$, et que pour $\nu$ variable, le système projectif des $(\underline\Phi_\nu^{(q)})$ \inserted{(ou $\underline\Psi_\nu^{(q)}$)} \deleted{\unsure{est} définit de façon} admet une limite projective dans la catégorie des faisceaux $\ell$-adiques constructibles. S'il en était ainsi, on déduirait par passage à la limite \inserted{projective}, à partir des suites spectrales \deleted{\unsure{(1.14)}} \inserted{(\deleted{\unsure{III}}1.14)} resp. \deleted{\unsure{III}} (2.5), des suites spectrales \deleted{de faisceaux} \inserted{\unsure{à composantes cohomologiques|à composantes cohomologique}} $\ell$-adiques, \unsure{\S} C'est ce \unsure{qu'on|qui} \unsure{constatera|constate} dans un cas particulier typique au n° \unsure{5}, -- et il est plausible que le cas général puisse s'y ramener par les \uline{arguments} \unsure{standart|standard}, via la résolution des \uline{singularités}.
